\par\bigskip\pdfbookmark[0]{De nexu. Fragmentum III}{nexus-fragment-III}\centerline{[III.]}\medskip
Es sei $C$ der Complexus der Repräsentanten sämmtlicher Classen der formae proprie primitivae für den Determinant $-D$. Wir bezeichnen durch $(n)$ die Anzahl aller Darstellungen der Zahl $n$ durch Formen aus dem Complexus $C$. Es sei $p$ eine ungerade Primzahl. Dann ist

\SourceLeaf{293}
\[
\begin{aligned}
&1.\quad(pn)=(n)\quad\text{wenn }p\text{ ein Divisor von }D\\
&\left.\begin{aligned}2.\quad(pn)&=(n)+(h)\\3.\quad(pn)&=-(n)+(h)\end{aligned}\right\}
\ \text{wenn }p\text{ Nichtdivisor von }D\\
&\hspace{7em}\left\{\begin{aligned}&\text{Divisor von }xx+D\\&\text{Nichtdivisor von }xx+D\end{aligned}\right.
\end{aligned}
\]
wo $n=hp^\mu$, $\mu$ beliebig und $h$ nicht durch $p$ theilbar.
\[
\begin{aligned}
\text{Im Fall }1.\ \text{ist}\quad&(h)=(ph)=(pph)=(p^3h)\text{ etc.}\\
2.\quad&(ph)=2(h),\ (pph)=3(h),\ (p^3h)=4(h)\text{ etc.}\\
3.\quad&(ph)=0,\ (pph)=(h),\ (p^3h)=0,\ (p^4h)=(h)\text{ etc.}
\end{aligned}
\]
\begin{center}$\ast\qquad\ast$\par$\ast$\end{center}

Aus jeder Classis pr. pr. für den Determinans $=-D$, deren Anzahl $=\lambda$, sei eine Form ausgewählt, und der Complexus dieser Form sei $L$.

Man bezeichne durch $fA$ die Anzahl sämmtlicher Darstellungen der Zahl $A$ durch Formen aus $L$.

Es sei ferner $f(A;p)=f\frac A{p^\alpha}$, wenn $p^\alpha$ die höchste Potenz der Primzahl $p$ ist, welche $A$ misst; ferner $f(A;p,q)=f\frac A{p^\alpha q^\beta}$, wenn $q$ eine andere Primzahl deren höchste $A$ messende Potenz $=q^\beta$ und so ferner $f(A;p,q,r)=f\frac A{p^\alpha q^\beta r^\gamma}$ wenn $r$ eine dritte Primzahl, deren höchste Potenz $A$ messend $r^\gamma$ ist u. s. w.
\BookRule{.30\textwidth}
